\section{Attribution and discussion}
The primary experiments establish three facts. First, detector operating point strongly affects downstream MOT: changing resolution, NMS, and confidence changes the observation set used for association. Second, the frozen policy improves the declared high-resolution baseline on two tracker hosts in modern VisDrone evaluation, with positive paired HOTA intervals. Third, operating-point choices shape compute: v3 reduces measured detector time while adding a small controller overhead. The harder question is what part of that result should be called adaptation.

The matched static control uses resolution 1088, NMS IoU 0.45, and detector confidence 0.40. Its archived HOTA is approximately 44.530 for ByteTrack and 48.269 for OATrack, compared with v3 values of approximately 44.523 and 48.095. These values are close. Thus the comparison against \texttt{r1536\_n70} demonstrates the value of moving away from that declared baseline, but it does not isolate a consistent additional accuracy benefit from changing the operating point frame by frame.

The timing control points in the same direction. On the development split, the causal schedule has mean $\Delta$HOTA 2.446957 against the baseline, while a same-action-distribution shuffled schedule has 2.617228. The shuffled schedule is not evidence that random timing is globally superior; the two values are calibration statistics from one frozen control. It is evidence that the present experiment cannot claim that exact scene-conditioned timing is necessary for the baseline-relative gain.

This attribution result does not invalidate AC--MOT. A runtime controller can remain useful when much of its accuracy advantage is explained by choosing a good operating-point mixture. The controller still provides online compute allocation, a mechanism to respond to changing resource constraints, and a deployment interface in which operating points can be selected without retraining detector weights. In an application with a hard average budget, a policy that spends high-resolution calls selectively may be preferable to a uniformly expensive fixed point even when balanced HOTA is similar.

The interpretation also explains the tracker dependence. ByteTrack starts from a weaker modern baseline and shows a larger MOTA and HOTA change; OATrack starts stronger and shows a smaller VisDrone HOTA increment and an essentially flat UAVDT HOTA transfer. A detector-side policy changes the input to each association algorithm, but the way each tracker converts extra or missing observations into identities is different. It would therefore be misleading to report one pooled ``AC--MOT gain.''

The modern negative searches are part of this discussion. Resolution-only, NMS-only, cue-subset, and earlier joint formulations did not pass the predeclared robustness gate against a matched static point. The final v3 candidate was selected against \texttt{r1536\_n70}, not because the archive showed that scene timing beat every fixed point. The selection audit also records that Trial 7 was not the literal argmax of the documented rule under either eligibility reading. The freeze is preserved exactly and the deviation is reported as provenance.

The strongest scientifically supported conclusion is therefore: AC--MOT is an auditable training-free mechanism for shaping the detector observation stream and computational operating point of aerial MOT. It can improve a declared baseline and can transfer with tracker-dependent outcomes. The evidence does not yet establish that exact causal scene-conditioned switching timing adds a separate, consistent accuracy benefit beyond operating-point calibration and action-mixture effects. Future work can test that question explicitly with cost-matched randomized policies, more detectors and trackers, larger unseen datasets, and a predeclared dynamic-scheduling objective.

\subsection{Why the controller remains useful when timing is not isolated}
The attribution result distinguishes scientific effect from systems value. A calibrated static point can match v3 HOTA on the reported hosts, but a static point must spend its chosen cost on every frame. A dynamic controller can move among operating points as a resource policy, react to changing scene statistics, and expose an interface for a future compute budget. For example, a deployment may impose an average detector-time budget rather than a single fixed resolution; a controller can then be evaluated against cost-matched schedules while preserving the detector and tracker. The current freeze does not perform that new experiment, but it identifies the comparison that would be needed.

The controller also provides a diagnostic instrument. Because its state is explicit and its action tuples are finite, one can inspect which scenes consume high-resolution or low-confidence actions, compare state occupancy across tracker hosts, and identify when FP reduction is purchased with FN. This is different from claiming that the policy has learned the optimal action. The value is an auditable connection between scene evidence, detector operating point, downstream association, and measured latency.

The attribution boundary therefore improves rather than weakens the paper's contribution. A claim that every observed gain is caused by scene-conditioned switching would be contradicted by the matched-static and shuffled controls. A claim that detector operating-point selection is irrelevant would be contradicted by the baseline-relative changes and runtime decomposition. The frozen evidence supports the intermediate conclusion: operating-point calibration and compute shaping matter; the incremental accuracy value of exact causal timing remains an open, separately testable question.
